% Avsnitt 14.7, ur lectures/L14/README.md.

\section{Sammanfattning}
\label{c14:sec:review}

\needspace{6\baselineskip}
\subsection*{Det här ska du kunna nu}
\begin{itemize}
  \item Implementera ett skiftregister som serialiserar MSB först och sätter in stoppbitar enligt
    femregeln.
  \item Förklara varför stoppningsräknaren inte får nollställas när en ny grupp laddas.
  \item Skilja på \code{bit_valid} och \code{bit_done}, och säga vilken av dem CRC-motorn ska
    grindas på.
  \item Säga vad \code{stuff} är till för, och varför den behövs utanför modulen.
\end{itemize}

\needspace{6\baselineskip}
\subsection*{Frågor att testa dig själv med}
\begin{enumerate}
  \item Varför måste den första biten i en nyladdad grupp ligga ute direkt, i stället för vid nästa
    \code{shift}?
  \item Varför nollställs inte stoppningsräknaren vid omladdning?
  \item Vad skulle gå fel om \code{can_controller} matade CRC-motorn på \code{bit_done} i stället
    för på \code{bit_valid and not stuff}?
  \item En stoppbit har alltid motsatt värde mot de fem föregående. Varför räcker inte ”sätt in en
    nolla”?
  \item Varför väntar \code{done} en extra skiftning i stället för att gå hög direkt när sista
    databiten skiftats ut?
\end{enumerate}

\needspace{6\baselineskip}
\subsection*{Riktvärde}
Ungefär halvvägs genom projektet. Tre av fyra lövmoduler brukar vara klara här, och det är också
ungefär nu det blir tydligt vilken grupp som faktiskt granskar varandras kod och vilken som bara
klickar ”Approve”.

\needspace{6\baselineskip}
\subsection*{Blicken framåt}
\Chapref{c15:ch} bygger \code{rx_shift_reg}, mottagarsidans spegelbild: deserialisering, avstoppning
och flaggning av brott mot stoppbitsregeln. \Chapref{c16:ch} sätter sedan båda registren i arbete
inuti \code{can_controller}.
